\documentclass[12pt]{article}
\def\activityid{F09-U-v3}\usepackage[letterpaper,margin=.65in,top=.60in,bottom=.65in]{geometry}
\usepackage[T1]{fontenc}
\usepackage[default]{sourcesanspro}
\usepackage{microtype,tikz,array,tabularx,fancyhdr,amsmath,amssymb}
\usepackage[hidelinks]{hyperref}
\usetikzlibrary{calc,arrows.meta}
\setlength{\parindent}{0pt}\setlength{\parskip}{7pt}
\setlength{\headheight}{0pt}\setlength{\footskip}{26pt}\setlength{\emergencystretch}{2em}
\pagestyle{fancy}\fancyhf{}\renewcommand{\headrulewidth}{0pt}\renewcommand{\footrulewidth}{.4pt}
\fancyfoot[L]{\scriptsize Bellingham Math Circle / Week 9 / \activityid}
\fancyfoot[R]{\scriptsize \thepage}
\definecolor{ink}{gray}{.12}\definecolor{soft}{gray}{.95}\color{ink}
\newcommand{\PageTitle}[2]{{\fontsize{10}{12}\selectfont\bfseries WEEK 9\quad BOUNCE LAB\hfill #1\par}\vspace{3pt}{\fontsize{26}{29}\selectfont\bfseries #2\par}\vspace{2pt}}
\newcommand{\Subhead}[1]{\par\vspace{3pt}{\large\bfseries #1}\par}
\newcommand{\RuleBox}[1]{\par\begingroup\setlength{\fboxsep}{8pt}\colorbox{soft}{\parbox{\dimexpr\linewidth-16pt\relax}{#1}}\endgroup\par}
\newcommand{\WriteBox}[2]{\par\textbf{#1}\par\vspace{2pt}\begin{tikzpicture}\draw[black!40,line width=.5pt] (0,0) rectangle (\dimexpr\linewidth-.5pt\relax,#2);\end{tikzpicture}\par}
\newcommand{\RookBoard}[3]{\begin{tikzpicture}[x=#3,y=#3]\draw[step=1,black!45] (0,0) grid (#1,#2);\draw[thick] (0,0) rectangle (#1,#2);\node[font=\Large] at ({#1-.5},.5){$\star$};\draw[-{Latex},thick] (0,-.35)--(#1,-.35);\draw[-{Latex},thick] ({#1+.35},#2)--({#1+.35},0);\end{tikzpicture}}
\newcommand{\Table}[3]{\begin{tikzpicture}[x=#3,y=#3]\draw[step=1,black!25] (0,0) grid (#1,#2);\draw[thick] (0,0) rectangle (#1,#2);\fill (0,0)circle(3pt);\draw[-{Latex},very thick] (0,0)--(.7,.7);\node[below] at ({#1/2},-.1){width #1};\node[rotate=90,above] at (-.2,{#2/2}){height #2};\end{tikzpicture}}
\newcommand{\GraphNodes}{\coordinate(A)at(0,0);\coordinate(B)at(3,0);\coordinate(C)at(1.5,2.6);\coordinate(D)at(1.5,.95);}
\newcommand{\Kfour}[1]{\begin{tikzpicture}[scale=#1]\GraphNodes\draw[very thick](A)--(B)--(C)--(A)--(D)--(B) (D)--(C);\foreach \n in {A,B,C,D}{\node[circle,draw,fill=white,minimum size=22pt,font=\bfseries]at(\n){\n};}\end{tikzpicture}}
\newcommand{\House}[1]{\begin{tikzpicture}[scale=#1]\coordinate(A)at(0,0);\coordinate(B)at(3,0);\coordinate(C)at(3,2);\coordinate(D)at(0,2);\coordinate(E)at(1.5,3.3);\draw[very thick](A)--(B)--(C)--(D)--(A) (D)--(E)--(C);\foreach \n in {A,B,C,D,E}{\node[circle,draw,fill=white,minimum size=22pt,font=\bfseries]at(\n){\n};}\end{tikzpicture}}




\newcommand{\StudentPage}[1]{{\fontsize{10}{12}\selectfont\bfseries Week 9 / Billiards lab / #1\par}\vspace{10pt}}
\newcommand{\Problem}[2]{\par\noindent\textbf{Problem #1:} #2\par}
\newcommand{\Space}[1]{\par\begin{tikzpicture}\draw[black!35] (0,0) rectangle (\dimexpr\linewidth-.5pt\relax,#1);\end{tikzpicture}\par}
\newcommand{\Piles}[1]{\begin{center}\begin{tikzpicture}[x=1in,y=1in]\foreach \x in {0,...,#1}{\draw[black!50,rounded corners] ({\x*2.05},0) rectangle ({\x*2.05+1.8},1.3);}\end{tikzpicture}\end{center}}

\newcommand{\Rooms}[5]{\begin{tikzpicture}[x=#5,y=#5]\draw[step=1,black!20](0,0)grid(#3,#4);\foreach\x in{0,#1,...,#3}{\draw[thick](\x,0)--(\x,#4);}\foreach\y in{0,#2,...,#4}{\draw[thick](0,\y)--(#3,\y);}\fill(0,0)circle(3pt);\end{tikzpicture}}
\newcommand{\PlainGrid}[3]{\begin{tikzpicture}[x=#3,y=#3]\draw[step=1,black!30](0,0)grid(#1,#2);\fill(0,0)circle(3pt);\end{tikzpicture}}

\begin{document}

\StudentPage{Grades 4--5}
\Problem{1}{Start at each black dot with slope 1: one unit across for each unit up. Bounce at walls, reversing only the direction blocked by the wall. Stop at the first corner. Trace each path, mark wall bounces with dots, and circle the endpoint. Crossings of old path lines are not walls.}
\begin{center}\Table{2}{3}{.57in}\hspace{.8in}\Table{3}{2}{.57in}\end{center}
\begin{center}\Table{3}{3}{.51in}\end{center}
\Problem{2}{Record the ending corner and number of bounces for each table. Do not count the launch or final corner as bounces. Compare the two non-square tables.}
\begin{center}\renewcommand{\arraystretch}{1.8}\begin{tabular}{c|p{2in}|p{2in}}Width $\times$ height&Corner&Bounces\\\hline$2\times3$&&\\$3\times2$&&\\$3\times3$&&\end{tabular}\end{center}
\newpage
\StudentPage{Grades 4--5}
\Problem{3}{Trace the $4\times6$ table on the left. On the right are mirror copies of the same room. Draw a straight diagonal from the black dot to the first intersection of a vertical and horizontal room wall. Number the wall crossings before that intersection and match them to bounces on the left.}
\begin{center}\Table{4}{6}{.36in}\hspace{.7in}\Rooms{4}{6}{12}{12}{.26in}\end{center}
\Problem{4}{List the horizontal distances where the straight path crosses a vertical room wall. List the vertical distances where it crosses a horizontal room wall. Circle the first number appearing in both lists. The two distances agree along this diagonal.}
\Space{1.0in}
\Problem{5}{Use your circled distance to count whole room widths and heights. Move your finger back and forth across the original width for each whole width, then up and down for each whole height. Record the side reached in each direction and the resulting corner.}
\Space{1.0in}
\newpage
\StudentPage{Grades 4--5}
\Problem{6}{Use wall-distance lists to find the first common distance $L$ for each table. Divide it by the width to count whole widths, and by the height to count whole heights. Record the folded corner. Odd whole-width counts end right; even counts end left. Odd whole-height counts end top; even counts end bottom.}
\begin{center}\renewcommand{\arraystretch}{2.3}\begin{tabular}{c|c|c|c|p{1.25in}}Table&$L$&Widths&Heights&Corner\\\hline$4\times6$&&&&\\$6\times9$&&&&\\$6\times10$&&&&\\$8\times14$&&&&\end{tabular}\end{center}
\Problem{7}{Compare $4\times6$ with $6\times9$. Draw or record how the table and unfolded distances changed. Then compare $4\times6$ with $6\times10$, whose sides are also both even. Circle the two tables with different ending corners.}
\Space{1.2in}
\Problem{8}{On the $4\times6$ copied rooms, extend your imaginary diagonal to distance 24. It would cross 6 whole widths and 4 whole heights. Record the corner those counts predict, then mark the earlier corner where this game actually stops.}
\Space{.9in}
\newpage
\StudentPage{Grades 4--5}
\Problem{9}{For each table, list the distances of vertical and horizontal wall crossings before its first corner. Count all these crossings. Compare your counts with your earlier traced paths.}
\begin{center}\renewcommand{\arraystretch}{2.6}\begin{tabular}{c|p{1.7in}|p{1.7in}|c}Table&Vertical walls&Horizontal walls&Bounces\\\hline$2\times3$&&&\\$3\times3$&&&\\$4\times6$&&&\\$6\times10$&&&\end{tabular}\end{center}
\Problem{10}{Choose a table from Problem 9. Draw its copied rooms or record both wall lists. Circle its final corner separately. Count how many vertical crossings are left when the final corner is removed; then count horizontal crossings. Write a rule using whole-width count $m$ and whole-height count $n$.}
\Space{1.4in}
\Problem{11}{Make a scaled copy of your chosen rectangle by doubling both sides. Predict its corner and bounce count. Check using wall distances and record your result.}
\Space{1in}
\newpage
\StudentPage{Grades 4--5}
\Problem{12}{Try these three rectangles. Trace their paths and record corners and bounce counts. All have width plus height equal to 6. Circle those reaching top-right after exactly four bounces.}
\begin{center}\Table{1}{5}{.44in}\hspace{.45in}\Table{3}{3}{.44in}\hspace{.45in}\Table{5}{1}{.44in}\end{center}
\begin{center}\renewcommand{\arraystretch}{1.8}\begin{tabular}{c|p{2in}|p{2in}}Table&Corner&Bounces\\\hline$1\times5$&&\\$3\times3$&&\\$5\times1$&&\end{tabular}\end{center}
\Problem{13}{Make a different-size rectangle ending top-right after four bounces. Record its size, first common distance, and whole-room counts.}
\Space{.8in}
\Problem{14}{Try to make a table ending top-right after three bounces. Use the whole-room counts to list your candidates. Record which corner each candidate reaches, or which earlier corner stops the path.}
\Space{.9in}

\end{document}
